\chapter{Implementation: challenges and solutions}
\label{cha:implementation}

Turning the design of \Cref{cha:design} into a working, evaluable system surfaced a number of problems that were not visible at the design stage. Almost every difficulty turned out to be a case of a component quietly producing something that looked right, and almost every solution consisted of removing the opportunity to be quietly wrong rather than adding an instruction not to be.

\section{Keeping nutrients out of the model's output}
\label{sec:impl-portion-ref}

The first design considered for the agent's output schema included a nutrient breakdown alongside every portion, on the reasoning that having the model state what it believed each ingredient contributed would make its intermediate reasoning inspectable. This was abandoned almost immediately: a schema field that accepts a nutrient value is a schema field the model can, and eventually will, fill with an invented number -- exactly the hallucination risk described in \Cref{sec:background-llm-failure-modes}. The schema was cut down to the bare reference of \Cref{sec:design-no-invented-foods}, shown against the first design in \Cref{fig:diagram-portion-schema}. The cost of this decision is that the model can no longer show its own arithmetic; the benefit is that there is no field left in which it could hallucinate a number that later gets treated as ground truth.

\begin{figure}[htbp]
    \centering
    \resizebox{0.85\textwidth}{!}{\begin{tikzpicture}[
    rec/.style={draw, rectangle split, align=left, text width=42mm, inner sep=1.8mm, font=\small},
    title/.style={font=\small\bfseries, align=center, text width=46mm},
    strike/.style={draw=black!45, line width=0.6pt},
]

\node[rec, rectangle split parts=8] (before) {
    \textbf{portion}
    \nodepart{two} code
    \nodepart{three} name
    \nodepart{four} grams
    \nodepart{five} \textcolor{black!40}{energy\_kcal}
    \nodepart{six} \textcolor{black!40}{protein\_g}
    \nodepart{seven} \textcolor{black!40}{carbs\_g}
    \nodepart{eight} \textcolor{black!40}{fat\_g}
};
\node[title, above=3mm of before.north] {first design, rejected};

\draw[strike] (before.five west) -- (before.five east);
\draw[strike] (before.six west) -- (before.six east);
\draw[strike] (before.seven west) -- (before.seven east);
\draw[strike] (before.eight west) -- (before.eight east);

\node[rec, rectangle split parts=4, right=34mm of before.north east, anchor=north west] (after) {
    \textbf{portion}
    \nodepart{two} code
    \nodepart{three} name
    \nodepart{four} grams
};
\node[title, above=3mm of after.north] {adopted};

\draw[-{Latex[length=2mm]}] (before.four east)
    -- node[above, font=\footnotesize, align=center] {every nutrient\\ field deleted} (after.west |- before.four east);

\end{tikzpicture}
}
    \caption{The portion schema before and after. A field that accepts a nutrient value is a field a model can fill with an invented one, so every such field was removed rather than accompanied by an instruction not to use it.}
    \label{fig:diagram-portion-schema}
\end{figure}

\section{Dirty data on one side, real zeros on the other}
\label{sec:impl-missing-nutrients}

The two data sources behave so differently that the same absent nutrient field means opposite things depending on which table a food came from, and getting that right was the substantive data-handling work of the build.

Because Open Food Facts is crowd-sourced (\Cref{sec:tools-food-data}), its raw export cannot be used as-is: a build-time pass in the setup package cleans and filters it before it becomes the catalogue the runtime queries. A row is admitted only if it carries a per-\qty{100}{\gram} value for all four macronutrients -- energy, protein, carbohydrate and fat -- and is dropped outright if it has no usable identity at all, meaning no product name in either the English or Polish field and no category. Every nutrient amount is rescaled into one canonical unit, since contributors record the same nutrient in different units. An energy sanity pass then repairs the single most common transcription error, a decimal point off by a power of ten, by testing each candidate rescaling of the stated kilocalorie figure against an Atwater estimate computed from the row's own macronutrients and keeping the one that reconciles; it deletes any row still above a physically implausible ceiling, and any row whose stated energy contradicts its macronutrient-derived estimate beyond a tolerance band. The catalogue is additionally scoped to products marked as available in Poland. Even after all this, every nutrient beyond the four mandatory macros is optional and simply stays empty when no contributor ever transcribed it -- so an absent fibre value on an Open Food Facts food genuinely means ``unknown''.

USDA FoodData Central is the opposite kind of artefact: a laboratory-curated reference dataset in which a food's record lists the nutrients that were actually assayed. An absent nutrient there does not mean a contributor skipped a field; it means the food does not meaningfully contain that nutrient, so it is read as a real zero rather than as a gap. The system therefore interprets a missing value by source -- unknown for Open Food Facts, zero for USDA -- with the sole exception of a nutrient USDA does not model as a column at all (salt, which it reports only as sodium), which stays unknown because there is no assayed value to stand in for it either way.

This distinction is exactly what the aggregation layer must not flatten. Treating every absent field as zero when summing -- the naive strategy -- was rejected quickly, because for an Open Food Facts food it produces a total that looks exactly like a real, complete measurement while silently understating it: a plan whose fibre total reads \qty{8}{\gram} because half its ingredients simply have no fibre data on file is indistinguishable, from the number alone, from a plan that genuinely contains \qty{8}{\gram} of fibre. That is precisely the kind of confidently wrong number the totaller exists to prevent. The totaller instead tracks, per nutrient, which portions of the plan (by ingredient, and by the share of the plan's total mass they represent) carry no recorded value for that nutrient, and reports the resulting total together with a coverage warning whenever any contributing food is missing data, naming up to three of the largest offending foods by mass so that a downstream reader -- either a human or, more importantly, the critic in the reflection loop -- knows which total to distrust and by how much. Because a USDA absence now resolves to a genuine zero rather than to a gap, these warnings in practice fire almost only on the crowd-sourced Open Food Facts half of the catalogue, which is where the real uncertainty lives. \Cref{alg:coverage} sketches the computation.

\begin{algorithm}[htbp]
\caption{Coverage-aware nutrient total with per-nutrient warnings}
\label{alg:coverage}
\begin{algorithmic}[1]
\Function{TotalWithWarnings}{$portions$}
    \State $totals \gets$ empty map from nutrient to exact rational sum
    \State $missingMass \gets$ empty map from nutrient to (total mass, list of foods)
    \For{each $portion$ in $portions$}
        \For{each $nutrient$ tracked}
            \If{$portion.food$ has a value for $nutrient$}
                \State $totals[nutrient] \gets totals[nutrient] + portion.food[nutrient] \times \frac{portion.grams}{100}$
            \Else
                \State record $portion.grams$ and $portion.food.name$ under $missingMass[nutrient]$
            \EndIf
        \EndFor
    \EndFor
    \State $warnings \gets$ empty list
    \For{each $nutrient$ with a nonempty $missingMass[nutrient]$}
        \State $\mathit{share} \gets$ missing mass as a fraction of the plan's total mass
        \State append a warning naming $nutrient$, $\mathit{share}$, and up to three foods from $missingMass[nutrient]$
    \EndFor
    \State \Return $(totals, warnings)$
\EndFunction
\end{algorithmic}
\end{algorithm}

Because the totaller uses the exact rational arithmetic of \Cref{sec:tools-aggregation}, a reported total differs from a genuinely complete measurement only by the contribution of the foods listed as missing, never by an additional rounding error compounding the understatement.

\section{Choosing what text to embed per entry}
\label{sec:impl-hybrid-search}

The hybrid retrieval strategy itself -- a lexical channel, a dense channel, and reciprocal rank fusion over the two -- was settled in \Cref{sec:tools-retrieval} and wired into the lookup tool in \Cref{sec:design-hybrid-search}. What that design leaves open, and what building the index actually forced a decision on, is exactly which text is fed to the embedding model for each catalogue entry, since the dense channel is only as good as the string it embeds. Embedding only the product name discarded useful signal that a name alone often does not carry (a category, a brand); embedding the entire raw record, on the other hand, diluted the resulting vector with data-quality noise specific to a crowd-sourced dataset (repeated placeholder text, inconsistent formatting, occasionally entire other languages left untranslated in a description field). The document embedded per Open Food Facts entry was settled on instead as a short, curated prose sentence assembled from the fields most likely to carry a food's identity -- its English name, its Polish name where that differs, a category phrase, and the brand -- so that a mozzarella product reads as ``Mozzarella (Mozzarella light). A mozzarella product by Galbani.'' rather than as a raw dump of every column. The USDA document is the analogous concatenation of a food's description and its category. The crowd-sourced ingredient list, though shown to the model in a truncated form on a lookup hit, is deliberately left out of the embedded text, since it is the single noisiest field in the record and the one most likely to pull an entry's vector off toward an unrelated cluster.

\section{Tool ergonomics under a limited context budget}
\label{sec:impl-tool-ergonomics}

An early version of the lookup tool accepted one query per call, mirroring how a human might type one search into a search box at a time. What makes that costly is the shape of a single tool call, drawn in \Cref{fig:diagram-tool-call-loop}: the runner sends the entire conversation so far to the model provider over the network, the provider replies with a tool call and generation halts, the runner executes the lookup locally against DuckDB, and then the runner must send the entire conversation \emph{again} -- now longer by the tool call and its result -- back over the network for generation to resume. The two provider hops are the expensive part, and each one re-uploads the whole, growing conversation; the database hop itself is local and negligible by comparison. Looking up ingredients one at a time therefore pays for a full pair of these network round trips per ingredient, and composing a single day's meals needs dozens of them, each also spending context-window space on the conversational scaffolding a tool call and its result carry.

\begin{figure}[htbp]
    \centering
    \resizebox{\textwidth}{!}{\begin{tikzpicture}[
    x=1mm, y=1mm,
    head/.style={draw, rounded corners, anchor=center, text width=34mm,
                 minimum height=9mm, inner sep=1.2mm, font=\scriptsize, align=center},
    life/.style={draw=black!45, dashed},
    net/.style={-{Latex[length=2mm]}},
    netret/.style={-{Latex[length=1.8mm]}, dashed},
    loc/.style={-{Latex[length=2mm]}, draw=black!55},
    locret/.style={-{Latex[length=1.8mm]}, dashed, draw=black!55},
    msg/.style={font=\scriptsize, align=center},
    tag/.style={font=\scriptsize\itshape, text=black!60, align=center},
    halt/.style={circle, fill=black, draw=black, inner sep=0pt, minimum size=1.6mm},
]

% --- lifelines ---
\node[head] (runner) at (0, 0) {Agent runner\\(local process)};
\node[head] (api) at (75, 0) {Model provider\\API};
\node[head] (kb) at (150, 0) {Food knowledge base\\(local DuckDB)};

\draw[life] (runner.south) -- (0, -86);
\draw[life] (api.south) -- (75, -86);
\draw[life] (kb.south) -- (150, -86);

% --- network round trip out to the provider ---
\draw[net] (0, -12) -- node[msg, above] {send the whole conversation} (75, -12);
\draw[netret] (75, -22) -- node[msg, above] {tool call returned} (0, -22);
\node[halt] at (75, -22) {};
\node[tag, anchor=west] at (78, -22) {halts for tool call};

% --- local hop to the database ---
\draw[loc] (0, -38) -- node[msg, above] {\texttt{lookup\_foods}(batch of queries)} (150, -38);
\draw[locret] (150, -46) -- node[msg, above] {ranked hits} (0, -46);

% --- second network round trip to resume ---
\draw[net] (0, -60) -- node[msg, above] {resend the whole conversation,\\ now plus the tool result} (75, -60);
\draw[netret] (75, -70) -- node[msg, above] {generation resumes} (0, -70);

\end{tikzpicture}
}
    \caption{One food-lookup tool call. Generation halts at the provider, and resuming it costs a second network round trip that re-sends the entire, now-longer conversation; the local database hop between them is comparatively free. Looking up $N$ ingredients one at a time multiplies the two network hops by $N$.}
    \label{fig:diagram-tool-call-loop}
\end{figure}

The obvious remedy is to collapse that fan-out into a single call, so the tool was changed to accept a batch of queries at once -- up to fifty per call -- so that all of a recipe's ingredients, or in the more aggressive case actually used all of a meal's, resolve in one round trip. Because the tool interface only permits this rather than compelling it, the same intent was reinforced through the prompt: the nutrition agent is instructed to ``search every ingredient in one batched \texttt{lookup\_foods} call'' and, wherever a follow-up search is unavoidable, to ``minimize the number of follow-up calls'' by running them in bulk, so the model is actively steered away from the one-query-at-a-time behaviour the round-trip cost penalises. A second, related change concerned the shape of the tool's response: an early version returned a natural-language summary of the top matches (``the best match for `mozzarella' is a 250g pack of...''), which read well but was harder for the model to parse reliably into exact codes and gram-scaled nutrient values than a compact, uniformly structured JSON result carrying the same information field by field; the response format was switched to structured JSON for this reason. Finally, because a lookup batch can, in principle, request an unbounded number of results, the number of hits returned per query was capped -- a small default of two hits per source, and a hard ceiling above it -- trading a small risk of missing a marginally better match further down the ranked list for a predictable, bounded contribution to the context budget every lookup call makes.

\section{Tuning the critique-and-refine loop}
\label{sec:impl-critic-design}

The loop's division of labour -- a tool-less, read-only critic that only lists issues, and a separate refiner that acts on them -- was fixed in the design (\Cref{sec:tools-self-correction,sec:design-reflection}); what building it left open was how many refinement passes to allow before giving up, and what to do if the critic and refiner have not converged by then. An unbounded loop was rejected outright, since nothing guarantees a critic and refiner pair will ever agree, and an unbounded retry budget is incompatible with \ref{req:interactive-runtime}'s requirement that a session complete in interactive time. A small, fixed iteration cap of three passes was chosen instead, with the loop falling back to returning the last refined plan, rather than failing the session, if the cap is reached with issues still outstanding. Three was settled on for two practical reasons rather than derived from anything. Each additional pass adds a full critic and refiner round trip, and reflection already roughly triples a run's wall-clock time and its token spend (\Cref{sec:results-overhead}), so a larger cap would have put a session outside \ref{req:interactive-runtime}'s interactive-time requirement and multiplied the paid model usage an entire evaluation grid consumes, which this project's budget did not stretch to. Within that limit, reading generated plans suggested three passes are usually enough: the objections a refiner can act on tend to be raised and resolved in the early passes.

A second, open question -- whether the critic should use a different, stronger model than the one generating the plan, on the reasoning that a reviewer benefits from more capability than the party being reviewed -- was investigated only partially. Using one model for generation, feedback and revision alike is the configuration the \textsc{Self-Refine} result this loop builds on was itself obtained under -- that work ``uses a single LLM as the generator, refiner and the feedback provider''~\cite{madaan2023selfrefine} -- so an equal-capability critic is the natural baseline to measure first rather than an arbitrary restriction. The pipeline does support configuring a different model per role, but the ablation study in \Cref{cha:methodology} holds the critic and the generator on the same model throughout, to keep the number of independent variables in the experimental grid manageable; \Cref{sec:conclusions-future-work} returns to a genuinely stronger reviewer as a direction for further work.

\section{Prompt-level fixes discovered through iteration}
\label{sec:impl-prompt-fixes}

The prompt architecture these fixes act on -- its two tiers, its fixed block skeleton, and the design principles behind them -- is set out in \Cref{sec:design-prompts}; what follows are the individual changes that reading actual generated plans, rather than any a priori design consideration, forced on top of it. A number of smaller problems surfaced this way, and were each addressed with a targeted prompt change rather than a structural one.

Early plans for users with no explicitly stated micronutrient concern nonetheless sometimes combined several implicit risk factors -- a vegetarian diet pattern together with a woman of reproductive age, say, which jointly imply an elevated iron requirement neither factor implies on its own -- without the plan's rationale acknowledging it. The nutrition agent's prompt was extended to explicitly instruct it to reason about implicit micronutrient interactions arising from a combination of profile factors, rather than only from any single factor in isolation.

Related to this, early plans treated the free-text rationale field purely as an explanation of the choices already made, leaving no natural place to mention that a nutrient target could reasonably be met only partly through food and would benefit from supplementation (vitamin~B12 alongside metformin, for instance). The rationale field's prompt instructions were extended to explicitly invite this kind of supplementation note where relevant, rather than requiring the model to either silently omit the concern or force an unrealistic amount of a food into a recipe purely to hit a number on paper.

A third recurring pattern concerned portion sizes that did not correspond to how a food is actually sold or used: a recipe calling for 90~g out of a 100~g pack, or 110~g when only 100~g is available, both of which are mildly implausible instructions for an actual shopper. Explicitly instructing the model to prefer using a whole package, or a simple fraction of one, measurably reduced this pattern, though it remains a soft preference rather than an enforced constraint, since some ingredients genuinely are used in odd quantities.

A fourth issue was cosmetic but recurring, and specific to the Open Food Facts half of the catalogue, where an entry's name is whatever a volunteer contributor happened to type on the label. A hit can therefore arrive with no English name at all, only a Polish one, or with a brand string standing in where a product name should be; the reader falls back through exactly that chain -- English name, then Polish name, then brand, then the literal ``unknown'' -- so a food can reach the model under an unhelpful string, which then surfaces in a recipe. The prompt was extended to let the model compose its own clear, human-readable display name from the rest of a hit's fields when the supplied name is not usable, while still requiring the underlying code to be a real, resolvable database reference: the display name is the model's own wording, but the code beneath it must still resolve in the database and remains subject to the no-invented-foods contract in \Cref{sec:design-no-invented-foods}. The curated USDA descriptions do not have this problem, so the instruction is only ever exercised on Open Food Facts hits.

The empty-request case is handled by treating an empty query as its own explicit case and omitting the session-request section of the prompt entirely rather than inserting an empty one, since an omitted section reads to the model as ``there is no additional request'' more reliably than an empty one does. How the prompt is assembled per ablation variant -- so that an ablated agent never reads about a tool it was denied -- is part of the switch itself (\Cref{sec:design-variant-config}).

The mechanisms described so far are easiest to see working together on a single concrete run. \Cref{cha:worked-example} walks through one real session of the \emph{regular} scenario planned by \texttt{gemini-3.7-flash} with the totaller and reflection loop both active: the lookup response, the agent plan as emitted, the hydrated breakfast, the full day of meals, the day's totals against the frozen targets, the coverage caveats, and the patient-facing rationale.\label{sec:impl-worked-example} That one session cost eight model requests and about 175{,}000 tokens over roughly 100 seconds of wall-clock time, with the agent making a single \texttt{lookup\_foods} call for its whole day of ingredients -- the batching behaviour \Cref{sec:impl-tool-ergonomics} was designed to encourage. Those are the numbers for one run on one profile; characterising cost and accuracy across models, variants and scenarios is the job of the evaluation harness and \Cref{cha:results}.

